\section{Related Work}\label{sec:related}
GRPO provides an effective way to tune a model's weights. It compares rewards within a group of attempts at the same prompt. But it requires informative groups, which are unlikely to be obtained from a controller with little initial task ability. If all attempts receive the same reward, their group-relative advantages are zero~\cite{deepseekmath}. Dynamic sampling filters groups and collects additional samples to maintain an informative training batch~\cite{dapo}. Filtering cannot create successful behavior when all sampled attempts fail. In our setting, we also change the starting states and use dense feedback during initial skill acquisition.

SimpleVLA-RL applies group-relative updates and dynamic sampling to VLA rollouts. Its experiments also show why the initial policy matters: on five RoboTwin 2.0 tasks, RL increased average success from 7.3\% to 25.4\% after supervised fine-tuning on 100 demonstrations per task, while success remained zero without task-specific fine-tuning~\cite{simplevlarl}. We use separately practiced skills to obtain useful behavior before learning the complete instruction.

iRe-VLA alternates RL on an action head with supervised learning on expert demonstrations and successful online trajectories. During RL, it freezes the vision-language model; supervised learning then updates the model using both earlier and newly collected data~\cite{irevla}. We also alternate interaction and supervised learning, but our skill-recovery bank contains successes from earlier and recent policies on the target robot. The supervised update helps recover skills which degrade while the shared controller learns another skill.

VLA models rely on vision and proprioception, but successful behavior on familiar robots does not establish that the model can use these observations effectively on a new embodiment. PLD provides lightweight residual specialists over a frozen VLA, collects trajectories through policy-aligned residual interventions, and distills the resulting data into the generalist~\cite{pld}. This allows the generalist to learn from behavior generated by residual RL. We retain a residual controller on the target robot and use successful skill trajectories both for recovery and for learning one composed instruction.

We use self-imitation learning in the sense of learning from the policy's own past successful behavior. Our supervised action-regression loss differs from the actor-critic objective of the original Self-Imitation Learning method~\cite{oh2018}.
